
Machine learning research depends on well-documented datasets, yet documentation adoption remains incomplete. On HuggingFace, a major dataset repository with over 100,000 datasets, approximately 40\% of datasets have dataset cards, and many existing cards are incomplete. This gap persists despite established frameworks such as Datasheets for Datasets and Model Cards that specify what information should be documented.

The core challenge is reducing documentation burden. Prior work has focused on defining documentation standards—creating templates and frameworks that specify what to document—while automated metadata extraction systems can capture structural properties like file formats and distributions but cannot generate semantic information such as dataset motivation, collection process, or known limitations. This leaves a gap between what can be automatically extracted and what researchers must manually author.

This work investigates whether AI-powered suggestion generation can reduce documentation friction. The key observation is that documentation differs from code generation in error tolerance: minor imperfections in documentation phrasing carry low cost, while bugs in code can have significant consequences. This asymmetry may create more favorable conditions for AI acceptance in documentation tasks compared to code assistance, where reported acceptance rates range from 65-75\% for tools like GitHub Copilot.

We conducted a pilot study with 75 researchers using an AI documentation copilot that generates contextual suggestions for dataset documentation fields. The system uses few-shot prompting with Llama-3-8B-Instruct, analyzing dataset properties to generate relevant suggestions. Our primary finding is that researchers accepted 92.0\% of suggestions (median per-user rate), substantially exceeding our conservative 70\% threshold. This acceptance rate was consistent across dataset types and accompanied by substantial modification rates, suggesting active engagement rather than passive acceptance.

This pilot validates that researchers will interact with AI-generated documentation suggestions at high rates. However, it does not establish whether this acceptance translates to improved documentation completeness, quality, or reduced time burden. Those outcomes require controlled evaluation beyond the scope of this proof-of-concept. By demonstrating that the prerequisite condition—user adoption—is achievable, this work provides foundation for investigating the complete value chain from suggestion acceptance to documentation quality improvement.
